\documentclass[12pt]{article}
\usepackage[utf8]{inputenc}
\usepackage[T1]{fontenc}
\usepackage{hyperref}
\title{O'du Language Resources for K-12 Curriculum Development}
\author{Compiled by Claude Sonnet 5 (claude-sonnet-5) for Andr\'as Kornai}
\date{}
\begin{document}
\maketitle

\section{Introduction}
O'du (Ơ Đu, Iduh, Tay Hat; ISO 639-3 tyh) is a Khmuic (Austroasiatic) language of Nghệ An province in Vietnam and of Laos. It is nearly extinct. Only a few elderly people speak it, and most Ơ Đu now speak Thai. Vietnam's 2019 census counted 428 people.

\section{Primary Educational Resources}
None found.

\section{Linguistic and Reference Materials}
\begin{itemize}
\item Wikipedia, O'du language: \url{https://en.wikipedia.org/wiki/O%CA%BCdu_language}
\item Wikipedia, O Du people: \url{https://en.wikipedia.org/wiki/O_Du_people}
\item MPI-EVA numerals page for O'du: \url{https://lingweb.eva.mpg.de/channumerals/O'du.htm}
\end{itemize}

\section{Religious and Community Resources}
\begin{itemize}
\item SOAS Glocal page on the Ơ Đu community in Vietnam: \url{https://glocal.soas.ac.uk/o-du-community/}
\item Joshua Project, O-Du in Vietnam: \url{https://www.joshuaproject.net/index.php/people_groups/14118/VM}
\item Joshua Project, O-Du in Laos: \url{https://www.joshuaproject.net/index.php/people_groups/14118/LA}
\end{itemize}

\section{Digital and Audio Resources}
None found.

\section{Curriculum Development Context}
Little material exists. Revitalisation would depend on documenting what the last speakers know. The SOAS page is the most relevant community source found.
\end{document}